\section{Expanded account of the computational analysis}
\label{sec:report-methods}
This chapter explains the executable operations behind the final mathematical report. It is derived from the pinned scripts and their saved records. It does not introduce new fitted values, new trajectories or an unreported reconstruction of an unavailable failed attempt.

\subsection{Recovery of the actual production configuration}
The entry point \texttt{core.py::reference} loads the historical Task 40 validation helper, obtains its nested production objects and reads the frozen WT state from the separate resting state file. It checks the state digest and the six stimulated parameter object digests against the frozen inputs. The source manifest is compared with Git blob hashes for the relevant production files. The audit reports 150 manifest entries, with 149 matches and one instructional metadata difference in \texttt{AGENTS.md}; that difference is not described as a numerical model change.

The nested model comprises the selected whole cell model, the NKCC stimulation wrapper and the minimal NBC wrapper, with a separate AE4 source class wrapper when required. A clone used for equilibrium calculations replaces the input by constant full stimulation while retaining the selected kinetic parameters. The parameter audit distinguishes this active configuration from dataclass defaults, legacy fields and constructor reference tuples. This is why a table of defaults would not reproduce the calculation.

For a buffer parameter perturbation, the code adjusts an unused constructor reference chloride value to preserve the constructor's fixed charge consistency. The actual frozen conserved initial condition used in the inverse trajectories is not replaced by that reference tuple. Consequently, a changed buffer amount may change the initial algebraic pH even though the conserved state is held fixed. These two operations must not be conflated with a new genotype resting solve.
\sourcefile{evidence/task44/core.py::reference, clone; evidence/task44/output/source_audit.json; evidence/task43/audit.py}

\subsection{Invertible scaling and independent reconstruction}
The scaling object uses bath sodium for $C_0$, the frozen cell and lumen volumes for $V_0$ and $L_0$, and the initial fully stimulated WT NKCC cycle flux for $J_0$. The flux scale is a cycle rate, not a chloride rate. Distinct amount scales are used for the cell and lumen. The time scale is $t_0=C_0V_0/J_0$ and the water flux scale is $J_0/C_0$. The complete diagonal state scaling retains the regulatory fraction with scale one.

The independent index list is $(0,1,2,3,5,6,7,8,9,11,12)$ in the zero based production ordering. Expansion reconstructs intracellular and luminal alkalinity from the exact charge identities before the production evaluator is called. Thus the independent vector field still depends on carbonate speciation and electrical closure at each evaluation. A second, explicit dimensionless right hand side and a separately assembled thirteen state right hand side provide checks that are stronger than merely printing the same production output twice.

The recorded dimensional versus dimensionless comparison has a maximum scaled state discrepancy of approximately $7.524\times10^{-10}$ and a dimensional cumulative water recovery discrepancy of approximately $1.443\times10^{-15}$ pL. These are agreement measures at saved comparison points. They do not establish a singular perturbation approximation or a global integration error theorem. The exact verification additionally reconstructs all laws and source terms at 84 random valid states, using seed 45044 and twelve states per source class. Those finite checks supplement the proofs; they do not replace them.
\sourcefile{evidence/task44/core.py::Scaling; evidence/task44/verify_exact.py; evidence/task44/output/exact_verification.json; evidence/task44/output/rescaling_verification.json}

\subsection{Four different experiments, not one interchangeable simulation}
The first experiment reproduces the frozen finite production protocol. Its saved reference uses a resting observable at exactly $t=0$ and begins production integration at $10^{-6}$ s with the shared conserved state. The audited production solver is Radau with relative tolerance $10^{-7}$, amount tolerance $10^{-10}$ fmol, volume tolerance $10^{-12}$ pL, regulatory tolerance $10^{-10}$ and maximum step 2 s. The dose labels do not generate the prescribed calcium and beta inputs through a dose response model.

The second experiment is the separate constant full stimulus equilibrium problem. The code evaluates that vector field at an arbitrary fully stimulated time, conventionally 1 s; the \texttt{time\_s} field in an equilibrium table is not an equilibration time or an experimentally observed delay. Equilibria are solved by a nonlinear root method, not identified by declaring a finite trajectory endpoint stationary.

The third experiment is the local expression step about a solved equilibrium. It uses the linearised state and output maps and therefore describes a derivative response to a small parameter perturbation. The fourth is the Task 41 inverse family, integrated through 600 s with an additional accumulated outflow variable and separately examined under continued stimulation. A conclusion established for one of these experiments is not automatically a conclusion about another.

The parent integration adapter uses Radau, relative tolerance $10^{-8}$, absolute tolerance $10^{-11}$ and maximum step 2 s for its standard analysis integration; the dimensionless tolerances are scaled consistently. The inverse verification uses the tighter relative tolerance $2\times10^{-9}$ and maximum step 1 s, including an independent thirteen state integration with subsequent quadrature of dense outflow. These analysis settings are not presented as the frozen production tolerances.
\sourcefile{evidence/task43/output/numerical_controls.json; evidence/task44/core.py::integrate; evidence/task44/inverse_analysis.py; evidence/task44/inverse_verify.py}

\subsection{Equilibrium solution and physiological acceptance}
For each source class and expression setting, ten positive independent chemical and volume coordinates are represented by their logarithms. The regulatory equilibrium is set to one. The hybrid nonlinear root method uses \texttt{xtol=1e-10} and a limit of 1500 function evaluations. The returned candidate is expanded to all thirteen production coordinates, and all eleven independent equations are then checked. In the root helper, a maximum independent residual below $10^{-7}$ is the numerical acceptance condition. The records separately retain the solver's success flag, message and evaluation count. Residual based acceptance must not be silently described as a theorem of existence or as a complete diagnostic of every solver status.

C0 is continued through expression levels $1,0.9,0.75,0.5,0.25,0.1,0.05,0$. Other classes use the converged C0 WT seed for their WT solve and matching C0 expression seeds for the 5\% and null solves. This produces 26 class and expression cases. Since all seven null source terms vanish, seven entries solve the same null system. The count is not a count of 26 distinct equilibria of one fixed vector field.

The early inherited checkpoint retained only an aggregate statement about four unsuccessful continuation attempts. Their individual seeds, class labels and error messages are not available in that record. The completed analysis preserves that limitation and supplies a documented successful retry strategy. It does not invent an attempt history and does not reinterpret failure of a particular numerical search as nonexistence of a root.

The initial helper's compact physiological filter is supplemented by the later full inherited gate review. The complete review distinguishes positivity, finite diagnostics, chemical intervals, source bookkeeping, current residuals and architecture checks. Four WT class cases remain outside the physiological gates despite solved stable roots. The supplementary sensitivity verifier explicitly reapplies the complete inherited gates to all 384 nearby cases; a successful nonlinear solve alone is not the acceptance criterion.
\sourcefile{evidence/task44/core.py::stationary, physiology; evidence/task44/equilibrium_review.py; evidence/task44/verify_sensitivity_gates.py; evidence/task44/output/equilibrium_review.json}

\subsection{Jacobians, spectra and boundary differentiation}
The dynamic Jacobian is formed in the eleven scaled independent coordinates. The coordinate difference is the relative step multiplied by the larger of the coordinate magnitude and 0.1. The two principal steps are $10^{-5}$ and $5\times10^{-6}$. Time is then restored by division by $t_0$. Eigenvalues from the two calculations are matched individually, rather than compared only through their largest real parts. The saved maximum matched relative discrepancy is approximately $5.43\times10^{-8}$.

At a regulatory boundary, the physical side or a smooth extension must be used. A centred difference that crosses the production clamp at $r=1$ halves that Jacobian column and produces a spurious regulatory mode. The corrected code uses a second order inward derivative. The resulting regulatory eigenvalue is $-1/30$ per second. The chemical modes remain substantially slower. The same issue is relevant when interpreting the derivative at the null expression endpoint: the source uses an inward expression difference where appropriate.

Singular values and condition numbers refer to the declared scaled coordinates. They are numerical conditioning diagnostics and cannot be interpreted as coordinate independent measures of biological fragility. Negative real eigenvalues support local stability of the computed root under the stated smooth physical linearisation. They do not prove global attraction, forward invariance of the physiological set, or that a trajectory beginning at the frozen resting state never leaves a gate.

The parameter sensitivity calculation further checks decay estimates using fourth order differentiation of the ten coordinate chemical block at two steps, together with the exact regulatory eigenvalue. This uses the block structure while retaining the complete eleven mode spectrum. No chemical or volume state is discarded to obtain a faster numerical calculation.
\sourcefile{evidence/task44/core.py::jacobian; evidence/task44/equilibrium_review.py; evidence/task44/sensitivity_analysis.py::resolved_decay}

\subsection{Implicit sensitivities and what a decomposition means}
At a locally regular equilibrium the state sensitivity is obtained by solving the linear system $F_z z_\mu=-F_\mu$. The output derivative includes both direct parameter dependence and the dependence mediated by the solved state. The implementation validates the result against independently solved neighbouring roots, rather than assuming that a finite difference of the original residual verifies the complete equilibrium response.

For expression, the comparison uses steps $10^{-4}$ and $5\times10^{-5}$ at expression 1, 0.5, 0.05 and 0, with a second order inward rule at zero. Six quantities are checked: outflow, NKCC cycle flux, AE4 chloride flux, total pump cycles, NHE flux and intracellular pH. All eleven state derivatives are also checked. Part IV prints the four derivative sets and their checks, including the zero denominator qualification for any compensation ratio.

For the sixteen uncertain parameters, the independent neighbouring parameter points are $p\exp(\pm h)$ for $h=10^{-3}$ and $5\times10^{-4}$. WT and null roots are resolved independently at every perturbed point. Fixed state derivatives use the smaller steps $10^{-5}$ and $5\times10^{-6}$ to resolve cancellations. The output includes derivatives and elasticities for WT secretion and pH, the stationary null deficit, NBC and non NBC alkalinity shares, AE4 demand, the conditional NHE and no NBC capacities, compensation and both slow decay times. These are local partial perturbations without a new WT recalibration.

The original \texttt{coordinate\_attribution.csv} is a chain rule accounting by state coordinate. In the outflow law, explicit dependence occurs through lumen volume, so a coordinate attribution can put the output derivative into that coordinate while all transporters influence it indirectly. Such a table is not a unique mechanistic assignment of causal percentages to NKCC, pump and other pathways. Part IV retains the table and the modal projection data so this limitation is visible.

An equilibrium time parameter can have zero stationary sensitivity while affecting a time dependent response. Likewise, effective capacity and a full stimulation multiplier can enter as a product and have identical stationary sensitivities without being separately identifiable. The source explicitly reports both examples. None of the chosen finite difference steps is an experimentally justified parameter range.
\sourcefile{evidence/task44/integrated_context.tex; evidence/task44/sensitivity_results.tex; evidence/task44/output/sensitivity_extended.json; evidence/task44/output/coordinate_attribution.csv}

\subsection{Modal response and nonnormal amplification}
The modal review uses the right eigenvector matrix and its inverse left rows, with right vectors normalised to Euclidean length one. For an expression perturbation, it evaluates each input projection, secretion output projection, transfer residue and contribution to the stationary derivative. The propagator norm is sampled over the recorded time grid extending to 10000 s. Its maximum of approximately 1.37 is a sampled result in the chosen state metric, not a proved supremum or a physiological energy gain.

The WT stationary derivative is small partly because modal contributions of opposite signs cancel. The slowest WT mode contributes approximately $-5.5754\times10^{-5}$ pL/s per unit expression, while the net derivative is approximately $1.2853\times10^{-6}$. The linear derivative at 600 s remains approximately $1.0027\times10^{-5}$, so the local approach to the stationary response is not instantaneous. The 30 s regulatory mode has zero expression input projection in this experiment because its target does not depend on expression. This does not imply that perturbing the regulatory input would have the same response.
\sourcefile{evidence/task44/equilibrium_results.tex; evidence/task44/output/equilibrium_review.json; evidence/task44/output/local_step_response.csv}

\subsection{Inverse scan, refinement, perturbation and extended stimulation}
The inverse calculation retains the one prescribed conductance family and the common conserved WT initial state. The 17 point ordered recruitment scan is followed by bracketed root refinement. Both the adaptive integral and original sampled trapezoidal integral are retained. The latter includes integer seconds together with the distinct initial points at zero and $10^{-6}$ s. The selected historical point is reproduced for WT, 5\% expression and null before a new comparison crossing is discussed.

The refined crossings are checked with two solver settings and against independently integrated production states with quadrature of dense flow. For each of eight effective parameters, the implicit crossing derivative is compared with independently solved neighbouring crossings at two logarithmic perturbation sizes. The WT denominator uses the same parameter perturbation as the null numerator. This is necessary to interpret the derivative of a ratio rather than a change in the null numerator alone.

Physiological gates are checked at accepted solver endpoints and the original observation grid. This is a finite numerical monitoring procedure, not a proof of interval wide invariance between every sampled point. The extended selected and crossing null cases are separate constant stimulus diagnostics through 3600 s. Their first pH crossing times are refined numerically; the late roots are labelled mathematical continuations outside the inherited pH interval. The report preserves the original 600 s observables alongside the later failure.
\sourcefile{evidence/task44/inverse_analysis.py; evidence/task44/inverse_verify.py; evidence/task44/inverse_detail/protocol.json; evidence/task44/inverse_detail/independent_verification.json}

\subsection{What clean replay verifies}
The original recorded clean replay applies to numerical checkpoint \nolinkurl{f783785a863440df459e9c5530beba4c7eb59110}. It ran eighteen recorded commands after removing generated numerical outputs and inverse caches in a separate disposable checkout. The comparison covers 47 authoritative JSON/CSV files: 46 regenerated products and one unchanged historical record. It reports 331157 bit identical numeric values, with explicitly listed runtime and ancestor validated commit metadata exclusions. The optional 532 trial cache files were regenerated rather than reused. A separate supplement records exact agreement for 7567 additional numeric values and the complete 384 nearby gate checks.

The present report assembly copies these receipts verbatim and verifies their file hashes and internal references. It does not claim to have repeated that entire numerical replay in this report writing session. Report compilation, reference checks, data table coverage, preservation of numerical payloads and visual review are recorded separately from the original scientific verification. The distinction between these two verification levels is retained in the final build record.
\sourcefile{evidence/task44/output/clean_recomputation_verification.json; evidence/task44/output/clean_sensitivity_gate_verification.json; evidence/task44/recompute_pinned_checkpoint.py}
